\chapter{Conclusion and Future Works}

\section{Conclusion}

This thesis presents a broad set of research studies across two main themes: representation-learning-based log anomaly detection, and the development and evaluation of image encryption and secret image sharing techniques. These seven chapters address key issues in system security and confidential data management.

A role-aware representation learning framework for detecting anomalies in logs is presented in Chapter~1. This framework learns token roles without requiring manual annotation or parsing, thereby overcoming the limitation of treating all tokens uniformly in previous methods such as NeuralLog. The proposed framework achieves an F1-score of 0.98--0.99 on four public datasets.

Attacks on chaos-based image encryption algorithms are discussed in Chapters~2 and~3. In Chapter~2, the medical image encryption algorithm of Jain et al.\ is fully broken using no more than $n+1$ oracle queries, with exact plaintext reconstruction (maximum pixel difference of 0), despite strong reported entropy and NPCR/UACI values. In Chapter~3, the same line of work is extended to break the MMCBIE image encryption algorithm, which uses the Hénon map, cyclic key mixing, and a two-dimensional logistic map, with at most 256 oracle queries by exploiting inadequate inter-pixel and inter-channel diffusion.

Pooling-based dimension reduction for secret image sharing is presented in Chapter~4 through PoolLiteSIS, which reduces shadow-image storage by 75\%. Lossless reconstruction is confirmed using PSNR~$=\infty$ and SSIM~$= 1.0$, while a residual permutation scheme improves security.

The robustness of XOR visual secret sharing schemes against learning-based attacks is investigated in Chapter~5. It is demonstrated that a neural network adversary could reconstruct the original image nearly perfectly (SSIM = 0.9720, PSNR = 33.87~dB) using all shares. This indicates that learning-based methods are capable of learning the entire reconstruction function. Nevertheless, in the missing share scenario, the quality of reconstructed images degrades substantially.

The image-sharing method in Chapter~6 used autoencoders to compress images into latent space and then distribute the latent vectors into different shares. The PSNR was over 30~dB and the SSIM score was close to 1.0. There was about 98\% compression from the original image size. However, each share contains little information about the original image.

Optimization strategies for secret image sharing are discussed in Chapter~7, including binary encoding, digit decomposition, average pooling with residual mapping, and pooling-based XOR methods. The Boolean-based $(t,n)$ MSIS method provides deterministic and perfect recovery using any $t$ consecutive shares or more than $t$. The autoencoder model provides the highest compression ratio without degrading performance, thereby supporting the goal of creating smaller shadows.

Overall, these contributions advance both anomaly detection and image security while consistently highlighting the distinction between performance-based evaluation metrics and true security properties.

% ─────────────────────────────────────────────────────────────────────────────

\section{Future Works}

This thesis introduces several research directions that can further expand the proposed methodologies in anomaly detection and image security by improving adaptability, robustness, and theoretical foundations.

For the first research direction, in log-based anomaly detection, there is substantial scope to improve the proposed role-aware learning methodology to make the framework more adaptable and efficient. In particular, learning the optimal number of roles and automatically adapting to changing log formats without retraining could improve applicability and flexibility. Furthermore, future work should consider multimodal system data, such as metrics and execution traces, hierarchical roles for improved system-state representation, and cross-system pretraining.

The cryptanalytic findings call for greater rigor in the analysis and development of chaotic cryptosystems. Future research may apply the proposed cryptanalysis techniques to a wider variety of cryptographic structures, from multidimensional chaos generators to permutation-diffusion schemes. The identified vulnerabilities can guide the design of stronger encryption schemes by implementing robust diffusion, adding channel dependency, and preventing deterministic keystream construction. Developing effective evaluation frameworks that combine statistical testing with resistance to structural attacks will substantially improve security.

Pooling-based secret image sharing schemes can adopt adaptive algorithms that better balance compression ratio and reconstruction quality. The effectiveness of data-driven reconstruction attacks highlights the need for schemes that are resistant to such attacks, potentially through share obfuscation or adversarially robust training. Autoencoder-based methods could further benefit from perceptual losses and generative learning to improve reconstruction quality while retaining high compression ratios. Video data sharing is another promising direction, where temporal redundancy can be exploited. For threshold schemes, further reduction of share size while maintaining perfect reconstruction remains an important area for improvement.

In all contributions, one recurring issue is the mismatch between the empirically observed metrics and theoretical guarantees. The future directions of research must concentrate on designing methods for evaluating solutions in light of these two factors. Furthermore, implementation on constrained hardware, including edge computing and IoT platforms, necessitates more efficient algorithms.
